\section{Introduction}
\label{sec:introduction}

Geodesics of black-hole spacetimes are integrated numerically in two very different regimes.
Ray tracing for images of accretion flows, such as those behind the Event Horizon Telescope
results~\cite{eht2019}, follows millions of photons for a short affine time each; codes such
as GRay~\cite{chan2013}, geokerr~\cite{dexter2009}, RAPTOR~\cite{bronzwaer2018} and
GYOTO~\cite{vincent2011} are built around adaptive Runge--Kutta methods or semi-analytic
solutions. The evolution of bound orbits, as in the adiabatic modelling of
extreme-mass-ratio inspirals, follows a few orbits for $10^3$--$10^5$ revolutions, where the
long-time behaviour of the integration error matters more than its size after one step.
Geometric integrators~\cite{gni2006} are designed for the second regime, and several have been
adapted to geodesics: Gauss collocation with a symmetric step-size control~\cite{seyrich2012},
explicit symplectic splittings for Schwarzschild~\cite{wang2021}, and Tao's extended phase
space~\cite{tao2016,pihajoki2015}, which the FANTASY code applies to arbitrary
metrics~\cite{christian2021}.

Comparisons between these families are usually made on one test problem, one formulation and
one length of integration. This report compares them systematically, on the simplest
spacetime where every relevant quantity has a closed form. We ask five research questions:

\begin{description}
\item[\RQ{1}] How do the mass-shell error and the phase error grow over long integrations?
\item[\RQ{2}] Does the formulation, the second-order geodesic equation or Hamilton's
  equations, change which quantities are conserved, and is a symmetric method that is not
  symplectic in the chosen variables enough?
\item[\RQ{3}] Which method reaches a given accuracy at the lowest cost, and how does the answer
  depend on the length of the integration?
\item[\RQ{4}] How far can an integrator extend the life of an orbit on the unstable photon
  sphere, and of an orbit at the marginally stable ISCO?
\item[\RQ{5}] Which integrator should one use for ray tracing, and which for long orbit
  evolution?
\end{description}

The study has four design choices that the results depend on. First, both formulations keep
the full eight-dimensional phase space, so that the conservation of the energy $E$ and the
axial angular momentum $L_z$ can be tested rather than imposed. Second, the integrators see
only a vector field, so the same code runs on toy problems, on Schwarzschild and, later, on
Kerr. Third, the cost of a run is the number of vector-field evaluations, counting every
fixed-point iteration of the implicit methods. Fourth, errors are measured on physical
quantities against exact references: the deflection angle at the radius where the
integration stops, the angle swept per radial period and the orbit $r(\phi)$ through elliptic
functions, and the exit time from the photon sphere through an elliptic integral.

Section~\ref{sec:formulations} sets out the two formulations and their invariants,
Sec.~\ref{sec:integrators} the integrators and Sec.~\ref{sec:testcases} the test problems and
references. Section~\ref{sec:results} presents the results, ordered by research question.
Section~\ref{sec:discussion} turns them into recommendations, and Sec.~\ref{sec:conclusions}
concludes.
